% =========================================================================
% thesis-config.tex - the file you edit on day one
%
% Student Guide:
% - Section 1 holds your thesis title, names, IDs, supervisors and semester.
%   Change the text inside the SECOND pair of braces only.
% - Every page (cover, approval, declaration, headers, forms) reads these
%   values, so you type each fact only once.
% - Sections 2 to 9 control the layout. Leave them alone.
% =========================================================================

% -------------------------------------------------------------------------
% 1. Your details (edit this section)
% -------------------------------------------------------------------------
\newcommand{\thesisTitle}{Write Your Thesis Title Here And Replace All Instructional Text With Your Actual Thesis Content After Reading The Instructions.}
\newcommand{\thesisSubtitle}{A Thesis in the Partial Fulfillment of the Requirements\\ for the Award of Bachelor of Computer Science and Engineering (BCSE)}

% Authors.
% Working alone? Make \authorB and \authorBID empty, like this:
%   \newcommand{\authorB}{}
%   \newcommand{\authorBID}{}
% Appendix A then shows one student by itself. The pages in 'frontmatter/'
% need a few small edits: see 'Thesis-Writing-Manual.pdf', section
% "Writing a Single-Author Thesis".
\newcommand{\authorA}{First Author}
\newcommand{\authorAID}{Student ID 1}
\newcommand{\authorB}{Second Author}
\newcommand{\authorBID}{Student ID 2}

% \ifsecondauthor{<text for two authors>}{<text for one author>}
% (used by Appendix A; do not edit)
\usepackage{etoolbox}
\newcommand{\ifsecondauthor}[2]{\ifdefempty{\authorB}{#2}{#1}}

% Supervisors and department leaders
\newcommand{\supervisorName}{Supervisor Name}
\newcommand{\supervisorDesignation}{Supervisor and Designation of the Supervisor}
\newcommand{\coSupervisorName}{Dr. Md. Rashedul Islam}
\newcommand{\coSupervisorDesignation}{Co-supervisor, Coordinator and Associate Professor}
\newcommand{\deanChairName}{Prof. Dr. Utpal Kanti Das}
\newcommand{\deanChairDesignation}{Dean and Chair}

% Institution details
\newcommand{\deptName}{Department of Computer Science and Engineering}
\newcommand{\schoolName}{IUBAT School of Computer Science and Engineering}
\newcommand{\univName}{IUBAT---International University of Business Agriculture and Technology}
\newcommand{\univAddress}{4 Embankment Drive Road, Sector 10, Uttara Model Town, Dhaka 1230, Bangladesh}
\newcommand{\submissionSemester}{Spring 2026}
\newcommand{\courseCodeTitle}{CSC 4000: Write Course Code and Title Here}

% -------------------------------------------------------------------------
% 2. Page size and margins (IUBAT thesis layout)
% Left margin 1.25 in (room for binding); top, bottom and right 1 in.
% -------------------------------------------------------------------------
\usepackage[
  a4paper,
  left=1.25in,
  right=1.0in,
  top=1.0in,
  bottom=1.0in,
  headheight=0.25in,
  headsep=0.3in,
  footskip=0.5in
]{geometry}

% -------------------------------------------------------------------------
% 3. Text and line spacing
% The front pages use LaTeX's normal double spacing. The main body
% (Chapter 1 onwards) switches to \bodytextmetrics below, which copies the
% departmental Word template:
%   * 27.5 pt from one line to the next (Word's "Double" for 12 pt text)
%   * 21.5 pt indent at the start of each paragraph (about 0.3 in)
% LaTeX's own \doublespacing is only about 24 pt, so the stretch factor
% below is set by hand.
% -------------------------------------------------------------------------
\usepackage[T1]{fontenc}
\usepackage{lmodern}
\usepackage{setspace}
\doublespacing

\usepackage{indentfirst}
\setlength{\parindent}{0.5in}
\setlength{\parskip}{0pt}

% Sizes here are in 'bp' (1/72 in), the unit PDF files are measured in.
% LaTeX's own 'pt' is slightly smaller, so 27.5pt would measure only 27.4
% in the PDF, while 27.5bp measures exactly 27.5.
% To change the body spacing, edit these two numbers.
\newcommand{\bodyStretchFactor}{1.90366}  % 14.5pt baseline -> 27.50bp pitch
\newcommand{\bodyTextIndent}{21.5bp}

% Turns on the body spacing and the Word-style list layout (list text
% 0.5 in from the margin, roman numerals in front, 6 pt between numbered
% items, no gap between bullets). main.tex calls it once, where Chapter 1
% begins, so the front pages are not affected.
\newcommand{\bodytextmetrics}{%
  \setstretch{\bodyStretchFactor}%
  \setlength{\parindent}{\bodyTextIndent}%
  \setlist[enumerate,1]{label=\Roman*.,align=right,leftmargin=36bp,
                        labelwidth=28bp,labelsep=8bp,
                        topsep=0pt,partopsep=0pt,itemsep=0pt,parsep=6bp}%
  \setlist[itemize,1]{leftmargin=36bp,labelwidth=12bp,labelsep=6bp,align=left,
                      topsep=0pt,partopsep=0pt,itemsep=0pt,parsep=0pt}%
  \setlist[enumerate,2]{label=\arabic*.,leftmargin=24pt,topsep=0pt,itemsep=3pt}%
  \setlist[itemize,2]{leftmargin=24pt,topsep=0pt,itemsep=0pt}%
}

% -------------------------------------------------------------------------
% 3b. Fixed-position blocks (used by the cover and the approval page)
% Do not edit. Each block is a box of fixed height, so a longer or shorter
% thesis title never moves anything else on the page.
%   \fixedblock{height}{content}: one box. The height is a plain number in
%       bp and is the distance down to the top of the next block.
%   \firstlinestrut: gives a block's first line a fixed height, so its
%       baseline always sits at the same place.
% -------------------------------------------------------------------------
\newcommand{\firstlinestrut}{\rule{0pt}{14bp}}
\newcommand{\fixedblock}[2]{%
  \vbox to #1bp{\parskip=0pt \parindent=0pt \centering #2\vfil}%
  \nointerlineskip}
% The thesis title (18 pt bold). #1 = height reserved for it.
\newcommand{\fixedtitleblock}[1]{%
  \fixedblock{#1}{\vskip-1.74bp
    {\fontsize{18pt}{24pt}\selectfont\bfseries\noindent\firstlinestrut
     \thesisTitle\par}}}
% One signature: line, name, designation.
% #1 = height of the block, #2 = name, #3 = designation.
\newcommand{\fixedsignatureblock}[3]{%
  \fixedblock{#1}{\noindent\vtop{\hsize=11cm \parskip=0pt \parindent=0pt \centering
    \rule{6.5cm}{0.4pt}\par\nointerlineskip\vskip3.3bp
    {\fontsize{12pt}{16pt}\selectfont\bfseries\noindent\firstlinestrut #2\par}
    {\fontsize{12pt}{16pt}\selectfont\noindent #3\par}}\par}}

% Lets LaTeX stretch spaces a little, so long URLs and words
% stay inside the margins.
\setlength{\emergencystretch}{1.5em}

% -------------------------------------------------------------------------
% 4. Header and footer (page number centered at the bottom)
% -------------------------------------------------------------------------
\usepackage{fancyhdr}
\fancypagestyle{plain}{
  \fancyhf{}
  \renewcommand{\headrulewidth}{0pt}
  \renewcommand{\footrulewidth}{0pt}
  \fancyfoot[C]{\thepage}
}
\pagestyle{plain}

% -------------------------------------------------------------------------
% 5. Chapter and section headings
% Chapter titles: 14 pt bold, centered.
% Section titles: 12 pt bold, left-aligned.
% -------------------------------------------------------------------------
\usepackage{titlesec}
\usepackage[normalem]{ulem}

\titleformat{\chapter}[block]
  {\normalfont\fontsize{14pt}{18pt}\selectfont\bfseries\centering}
  {Chapter \thechapter.\ }{0pt}{}

\titleformat{name=\chapter,numberless}[block]
  {\normalfont\fontsize{14pt}{18pt}\selectfont\bfseries\centering}
  {}{0pt}{}

% Gaps copied from the Word template (baseline to baseline, in bp):
%   chapter title to first body line    : 54.5
%   last body line to next section title: 53.5
%   subsection titles                   : no extra space
\titlespacing*{\chapter}{0pt}{-8.6bp}{27bp}

\titleformat{\section}[block]
  {\normalfont\fontsize{12pt}{15pt}\selectfont\bfseries}
  {\thesection\hspace{0.5em}}{0pt}{}
\titlespacing*{\section}{0pt}{25.05bp}{0pt}

\titleformat{\subsection}[block]
  {\normalfont\fontsize{12pt}{15pt}\selectfont\bfseries}
  {\thesubsection\hspace{0.5em}}{0pt}{}
\titlespacing*{\subsection}{0pt}{0.05bp}{0pt}

\titleformat{\subsubsection}[block]
  {\normalfont\fontsize{12pt}{15pt}\selectfont\bfseries\itshape}
  {\thesubsubsection\hspace{0.5em}}{0pt}{}
\titlespacing*{\subsubsection}{0pt}{0.05bp}{0pt}

% By default the 'report' class numbers headings only down to \subsection.
% These two lines also number sub-subsections and list them in the Table
% of Contents.
\setcounter{secnumdepth}{3}
\setcounter{tocdepth}{3}

% -------------------------------------------------------------------------
% 6. Table of Contents, List of Figures, List of Tables
% Dotted lines join each title to its page number.
% -------------------------------------------------------------------------
\usepackage{tocloft}

\renewcommand{\contentsname}{Table of Contents}
\renewcommand{\cfttoctitlefont}{\hfill\normalfont\fontsize{14pt}{18pt}\selectfont\bfseries}
\renewcommand{\cftaftertoctitle}{\hfill}

\renewcommand{\listfigurename}{List of Figures}
\renewcommand{\cftloftitlefont}{\hfill\normalfont\fontsize{14pt}{18pt}\selectfont\bfseries}
\renewcommand{\cftafterloftitle}{\hfill}

\renewcommand{\listtablename}{List of Tables}
\renewcommand{\cftlottitlefont}{\hfill\normalfont\fontsize{14pt}{18pt}\selectfont\bfseries}
\renewcommand{\cftafterlottitle}{\hfill}

% Dotted lines for chapters
\renewcommand{\cftchapleader}{\cftdotfill{\cftdotsep}}
\renewcommand{\cftchappagefont}{\normalfont}
\renewcommand{\cftchapfont}{\normalfont\bfseries}

% Words printed before the numbers
\renewcommand{\cftchappresnum}{Chapter }
\renewcommand{\cftchapaftersnum}{.}
\setlength{\cftchapnumwidth}{5.6em}

\renewcommand{\cftfigpresnum}{Figure }
\renewcommand{\cftfigaftersnum}{}
\setlength{\cftfignumwidth}{5.2em}

\renewcommand{\cfttabpresnum}{Table }
\renewcommand{\cfttabaftersnum}{}
\setlength{\cfttabnumwidth}{5.0em}

\setlength{\cftbeforetoctitleskip}{18pt}
\setlength{\cftaftertoctitleskip}{18pt}
\setlength{\cftbeforeloftitleskip}{18pt}
\setlength{\cftafterloftitleskip}{18pt}
\setlength{\cftbeforelottitleskip}{18pt}
\setlength{\cftafterlottitleskip}{18pt}

% -------------------------------------------------------------------------
% 7. Captions for figures and tables
% Centered, full 12 pt size. Figure caption below the figure, table caption
% above the table.
% -------------------------------------------------------------------------
\usepackage{caption}
\usepackage{subcaption}

% As in the Word template: plain (not bold) label and one space after the
% number, e.g. "Figure 3.1 Example Figure Title." Figure captions use the
% body line spacing; table captions are single-spaced.
\captionsetup{
  font={normalsize},
  labelfont={normalfont},
  labelsep=space,
  justification=centering
}
% Figure caption: below the figure, 6 pt gap.
\captionsetup[figure]{font={normalsize,stretch=\bodyStretchFactor},skip=6pt}
% Table caption: above the table, 12 pt gap, single-spaced.
\captionsetup[table]{font={normalsize,stretch=1.0},skip=12pt}

% -------------------------------------------------------------------------
% 8. Table and graphics packages
% Pictures, tables (including tables that continue on the next page),
% lists and maths.
% -------------------------------------------------------------------------
\usepackage{graphicx}
\usepackage{booktabs}
\usepackage{tabularx}
\usepackage{longtable}
\usepackage{xltabular}
\usepackage{array}
\usepackage{multirow}
\usepackage{enumitem}
% The list layout comes from \bodytextmetrics (section 3), so a plain
% \begin{enumerate} or \begin{itemize} in a chapter needs no options.
\usepackage{pdfpages}
\usepackage{amsmath,amssymb}

% -------------------------------------------------------------------------
% 9. Citations and links (IEEE style)
% -------------------------------------------------------------------------
\usepackage{cite}
\renewcommand{\bibname}{References}
\usepackage[pdfpagelabels=true,plainpages=false]{hyperref}
\hypersetup{
  colorlinks=true,
  linkcolor=black,
  citecolor=black,
  urlcolor=black,
  pdfborder={0 0 0}
}
